\documentclass[conference]{IEEEtran}
\usepackage{amsmath,amssymb}
\usepackage{booktabs}
\usepackage{url}

\begin{document}

\title{Barren Plateaus: Global vs.\ Local Cost Functions\\
\large Initial Project Synopsis --- AI 631 Quantum Machine Learning, Fall 2026 (Project \#8, QNN track)}

\author{\IEEEauthorblockN{Saif ur Rehman}
% Add group members here once confirmed (groups of 2--3).
\IEEEauthorblockA{GitHub: \url{https://github.com/Saifix/Barren-plateaus-global-vs.-local-cost-functions}}}

\maketitle

\section{Research Question}
Parameterized quantum circuits (PQCs) trained by gradient descent can suffer
from \emph{barren plateaus}: the variance of the cost gradient over random
parameters vanishes exponentially with the number of qubits $n$, so gradients
become too small to estimate with a polynomial number of shots~\cite{mcclean2018,larocca2025}.
Cerezo \emph{et al.}~\cite{cerezo2021} showed that this depends on the cost
function: for alternating layered ansatzes, a \emph{global} cost has
exponentially vanishing gradients at any depth, while a \emph{local} cost has
at worst polynomially vanishing gradients when the depth is
$L=\mathcal{O}(\log n)$. The link between cost locality and barren plateaus
was also studied in~\cite{uvarov2021}. We ask:

\textbf{RQ.} How does the gradient variance scale with qubit count $n$ and depth
$L$ for global and local costs on a hardware-efficient ansatz, does a smarter
initialization (identity-block~\cite{grant2019} or reduced-domain) help, and how
do these effects show up when training a small classifier, in simulation, under
a noise model, and on IBM hardware?

\section{Background and Related Work}
For deep circuits, the gradient variance has been characterized exactly
through the circuit's dynamical Lie algebra~\cite{ragone2024,fontana2024}.
Zhang \emph{et al.}~\cite{zhang2024} gave a general lower bound on the gradient
variance for circuits built from local 2-designs. Letcher
\emph{et al.}~\cite{letcher2024} derived tighter gradient bounds without
$t$-design assumptions that can be estimated efficiently on classical hardware,
and argued that observables with a non-vanishing local term can avoid barren
plateaus. The same global/local split appears in quantum generative models:
the MMD loss acts as an observable that is low-bodied and trainable, or global
and untrainable, depending on the kernel~\cite{rudolph2024}.
Cerezo \emph{et al.}~\cite{cerezo2025} ask whether provable absence of barren
plateaus implies classical simulability.

\emph{Initialization.} Wang \emph{et al.}~\cite{wangqi2024} prove that shrinking
each parameter's initial domain in proportion to $1/\sqrt{L}$ makes the
gradient decay at most polynomially in $n$ and $L$. Classical initializers
(Xavier, He, LeCun, orthogonal) reduce gradient-variance decay compared with
uniform random initialization~\cite{kashif2024date}, and the range of the
initial distribution alone can mitigate barren plateaus~\cite{kashif2024icrc}.

\emph{Noise.} Unital noise causes noise-induced barren
plateaus~\cite{wang2021}, and this analysis has been extended to a class of
non-unital maps that includes amplitude damping~\cite{singkanipa2025}.
Mele \emph{et al.}~\cite{mele2026} show that noise truncates most circuits to
effectively logarithmic depth, and prove that under non-unital noise, costs
built from local observables do not have barren plateaus. Engineered
dissipation has also been proposed as a way to restore
trainability~\cite{sannia2024}.

\section{Hypotheses and Falsifiers}
\textbf{H1 (cost locality).} At $L=\lceil\log_2 n\rceil$, $\mathrm{Var}[\partial_k C_G]$
decays exponentially in $n$ while $\mathrm{Var}[\partial_k C_L]$ decays much more slowly.
\emph{Falsified if}, over $n=4,\dots,14$, the fitted exponential decay rate
$b$ (from $\log\mathrm{Var}=a-bn$) of the local cost is not smaller than that of
the global cost (overlapping 95\% bootstrap CIs), or if the two variances at
$n=14$ differ by less than a factor of 10.

\textbf{H2 (depth).} For the local cost, increasing $L$ from $\lceil\log_2 n\rceil$
to $L\propto n$ makes its variance decay faster with $n$.
\emph{Falsified if} the fitted decay rate does not increase with $L$.

\textbf{H3 (initialization).} Identity-block or reduced-domain initialization raises
the initial gradient variance of the global cost and speeds up training.
\emph{Falsified if}, at $n\ge 10$, neither scheme raises the initial gradient
variance by at least $10\times$ over uniform initialization, or neither reduces
the median number of iterations needed to reach a fixed target loss (10 seeds).

\textbf{H4 (learning).} In the classification task, a model trained through a
global readout learns more slowly than the same circuit trained through a
local readout, and the gap widens with $n$.
\emph{Falsified if} the test-accuracy curves (mean over seeds) are
indistinguishable within one standard deviation at the largest $n$ tested.
Negative results will be reported as such.

\section{Models and Costs}
\textbf{Ansatz.} Hardware-efficient layered circuit $V(\boldsymbol\theta)$ on $n$
qubits: each layer applies $R_Y(\theta)R_Z(\theta)$ to every qubit, then CZ
gates on pairs $(1,2),(3,4),\dots$ followed by $(2,3),(4,5),\dots$ (brickwork).
This gives $2nL$ parameters and $(n-1)L$ CZ gates. This ansatz is not
guaranteed to satisfy the assumptions in~\cite{cerezo2021}, so this is an
empirical test.

\textbf{Costs.} With $\rho=V(\boldsymbol\theta)\,|\psi_0\rangle\langle\psi_0|\,V(\boldsymbol\theta)^\dagger$,
\begin{align}
C_G &= 1-\langle 0^{\otimes n}|\rho|0^{\otimes n}\rangle,\\
C_L &= 1-\tfrac{1}{n}\textstyle\sum_{j=1}^{n}\langle 0|\rho_j|0\rangle,
\end{align}
where $\rho_j$ is the reduced state of qubit $j$, as defined in~\cite{cerezo2021}.
Both costs come from the same $Z$-basis measurement: $C_G$ uses the
all-zeros frequency and $C_L$ uses the single-qubit marginals.

\textbf{Sanity check.} For the product ansatz $\bigotimes_j R_X(\theta_j)$ on
$|0\rangle^{\otimes n}$ with $\theta_j\sim U[0,2\pi]$, a direct calculation gives
$\mathrm{Var}[\partial_1 C_G]=\tfrac18(\tfrac38)^{n-1}$ and
$\mathrm{Var}[\partial_1 C_L]=\tfrac{1}{8n^2}$. We will use this closed form as
a unit test for the variance-estimation code.

\textbf{Initializations.} (i) Uniform $\theta\sim U[0,2\pi]$;
(ii) reduced-domain $\theta\sim U[-a,a]$ with $a\propto 1/\sqrt{L}$, following~\cite{wangqi2024};
(iii) identity blocks~\cite{grant2019}, where each block is a random sub-circuit
followed by its adjoint, so the block starts as the identity.

\textbf{Training task.} A fixed random shallow circuit $W$ prepares
$|\psi_0\rangle=W|0\rangle$, and $V(\boldsymbol\theta)$ is trained to map it back
to $|0\rangle^{\otimes n}$. This keeps identity initialization from being
trivially optimal.

\textbf{Classifier.} Angle encoding of $n$ PCA features, followed by the same
ansatz. The class score is either global ($P(0^{\otimes n})$) or local
($\tfrac12(1+\tfrac1n\sum_j\langle Z_j\rangle)$), trained with binary
cross-entropy. Gradients are exact on the statevector simulator and use the
parameter-shift rule~\cite{schuld2019} with finite shots and on hardware.

\section{Data and Baselines}
Gradient-variance studies use random circuits only, so no external data is
needed. The classification task uses the Breast Cancer Wisconsin (Diagnostic)
dataset shipped with scikit-learn, reduced by PCA to $n$ features
($n\in\{4,6,8,10\}$). We use the course NB11 starter harness for stratified
folds and leakage-free preprocessing (scaler and PCA fit on the training fold
only), with fixed seeds. The classical baselines are majority class, logistic
regression, and an MLP with a parameter count matched to the quantum model. We
report MLP training curves on the same folds next to the quantum curves. For
gradient scaling, the reference is the analytic product-ansatz result above.

\section{Experimental Protocol and Resources}
\textbf{Noise-free sweeps:} $n\in\{2,4,\dots,14\}$, $L\in\{1,\lceil\log_2 n\rceil, n\}$,
200 parameter draws per $(n,L,\text{cost},\text{init})$. We report the variance of
$\partial_k C$ for a fixed $k$ and averaged over $k$, with bootstrap CIs.
\textbf{Noisy simulation:} Qiskit Aer with depolarizing (unital) and
amplitude-damping (non-unital) channels, plus a noise model taken from an IBM
fake backend. This separates cost-induced from noise-induced
plateaus~\cite{wang2021,singkanipa2025} and tests whether the local cost
holds up better under non-unital noise, as~\cite{mele2026} suggests.
\textbf{Hardware:} parameter-shift gradients on an IBM backend at
$n\in\{2,4,6\}$, $L=\lceil\log_2 n\rceil$, 30--50 random draws, 2 circuits per
gradient, 4096 shots, so at most 300 circuits in total, with both costs from
the same circuits. We report the shot-noise floor next to measured variances.
Software: PennyLane and Qiskit~\cite{qiskit2024}.

\section{Week-by-Week Plan}
{\small\noindent
\begin{tabular}{@{}p{0.17\columnwidth}p{0.75\columnwidth}@{}}
\toprule
Week (due) & Work \\
\midrule
0 (Nov 14) & Repo and environment; analytic unit test; ansatz and both costs; variance-estimation script. \\
1 (Nov 15) & Noise-free variance sweep for $n\le 14$, uniform initialization; classifier pipeline with NB11 folds; classical baselines; first quantum classifier. Report: introduction and methods. \\
2 (Nov 22) & Full $(n,L)$ sweeps; three initializations; training task; global vs.\ local classifier curves; noise-model runs and ablations. Report: results draft. \\
3 (Nov 29) & IBM hardware parameter-shift gradients at $n=2,4,6$; compare with the noise model; presentation if scheduled. \\
4 (Dec 6) & Final IEEE report (PDF and source), final talk, tagged GitHub release with saved results and scripts to reproduce every figure. \\
\bottomrule
\end{tabular}\par}

\begin{thebibliography}{20}
\bibitem{mcclean2018}
J.~R. McClean, S.~Boixo, V.~N. Smelyanskiy, R.~Babbush, and H.~Neven,
``Barren plateaus in quantum neural network training landscapes,''
\emph{Nat. Commun.}, vol.~9, art.~4812, 2018, doi: 10.1038/s41467-018-07090-4.

\bibitem{larocca2025}
M.~Larocca \emph{et al.},
``Barren plateaus in variational quantum computing,''
\emph{Nat. Rev. Phys.}, vol.~7, pp.~174--189, 2025, doi: 10.1038/s42254-025-00813-9.

\bibitem{cerezo2021}
M.~Cerezo, A.~Sone, T.~Volkoff, L.~Cincio, and P.~J. Coles,
``Cost function dependent barren plateaus in shallow parametrized quantum circuits,''
\emph{Nat. Commun.}, vol.~12, art.~1791, 2021, doi: 10.1038/s41467-021-21728-w.

\bibitem{uvarov2021}
A.~V. Uvarov and J.~D. Biamonte,
``On barren plateaus and cost function locality in variational quantum algorithms,''
\emph{J. Phys. A: Math. Theor.}, vol.~54, art.~245301, 2021, doi: 10.1088/1751-8121/abfac7.

\bibitem{grant2019}
E.~Grant, L.~Wossnig, M.~Ostaszewski, and M.~Benedetti,
``An initialization strategy for addressing barren plateaus in parametrized quantum circuits,''
\emph{Quantum}, vol.~3, p.~214, 2019, doi: 10.22331/q-2019-12-09-214.

\bibitem{ragone2024}
M.~Ragone \emph{et al.},
``A Lie algebraic theory of barren plateaus for deep parameterized quantum circuits,''
\emph{Nat. Commun.}, vol.~15, art.~7172, 2024, doi: 10.1038/s41467-024-49909-3.

\bibitem{fontana2024}
E.~Fontana \emph{et al.},
``Characterizing barren plateaus in quantum ans\"atze with the adjoint representation,''
\emph{Nat. Commun.}, vol.~15, art.~7171, 2024, doi: 10.1038/s41467-024-49910-w.

\bibitem{zhang2024}
H.-K. Zhang, S.~Liu, and S.-X. Zhang,
``Absence of barren plateaus in finite local-depth circuits with long-range entanglement,''
\emph{Phys. Rev. Lett.}, vol.~132, art.~150603, 2024, doi: 10.1103/PhysRevLett.132.150603.

\bibitem{letcher2024}
A.~Letcher, S.~Woerner, and C.~Zoufal,
``Tight and efficient gradient bounds for parameterized quantum circuits,''
\emph{Quantum}, vol.~8, p.~1484, 2024, doi: 10.22331/q-2024-09-25-1484.

\bibitem{rudolph2024}
M.~S. Rudolph \emph{et al.},
``Trainability barriers and opportunities in quantum generative modeling,''
\emph{npj Quantum Inf.}, vol.~10, art.~116, 2024, doi: 10.1038/s41534-024-00902-0.

\bibitem{cerezo2025}
M.~Cerezo \emph{et al.},
``Does provable absence of barren plateaus imply classical simulability?''
\emph{Nat. Commun.}, vol.~16, art.~7907, 2025, doi: 10.1038/s41467-025-63099-6.

\bibitem{wangqi2024}
Y.~Wang, B.~Qi, C.~Ferrie, and D.~Dong,
``Trainability enhancement of parameterized quantum circuits via reduced-domain parameter initialization,''
\emph{Phys. Rev. Appl.}, vol.~22, art.~054005, 2024, doi: 10.1103/PhysRevApplied.22.054005.

\bibitem{kashif2024date}
M.~Kashif, M.~Rashid, S.~Al-Kuwari, and M.~Shafique,
``Alleviating barren plateaus in parameterized quantum machine learning circuits: Investigating advanced parameter initialization strategies,''
in \emph{Proc. Design, Automation \& Test in Europe Conf. (DATE)}, 2024, pp.~1--6, doi: 10.23919/DATE58400.2024.10546644.

\bibitem{kashif2024icrc}
M.~Kashif and M.~Shafique,
``The dilemma of random parameter initialization and barren plateaus in variational quantum algorithms,''
in \emph{Proc. IEEE Int. Conf. Rebooting Computing (ICRC)}, 2024, pp.~1--8, doi: 10.1109/ICRC64395.2024.10937003.

\bibitem{wang2021}
S.~Wang, E.~Fontana, M.~Cerezo, K.~Sharma, A.~Sone, L.~Cincio, and P.~J. Coles,
``Noise-induced barren plateaus in variational quantum algorithms,''
\emph{Nat. Commun.}, vol.~12, art.~6961, 2021, doi: 10.1038/s41467-021-27045-6.

\bibitem{singkanipa2025}
P.~Singkanipa and D.~A. Lidar,
``Beyond unital noise in variational quantum algorithms: Noise-induced barren plateaus and limit sets,''
\emph{Quantum}, vol.~9, p.~1617, 2025, doi: 10.22331/q-2025-01-30-1617.

\bibitem{mele2026}
A.~A. Mele \emph{et al.},
``Noise-induced shallow circuits and the absence of barren plateaus,''
\emph{Nat. Phys.}, vol.~22, pp.~751--756, 2026, doi: 10.1038/s41567-026-03245-z.

\bibitem{sannia2024}
A.~Sannia, F.~Tacchino, I.~Tavernelli, G.~L. Giorgi, and R.~Zambrini,
``Engineered dissipation to mitigate barren plateaus,''
\emph{npj Quantum Inf.}, vol.~10, art.~81, 2024, doi: 10.1038/s41534-024-00875-0.

\bibitem{schuld2019}
M.~Schuld, V.~Bergholm, C.~Gogolin, J.~Izaac, and N.~Killoran,
``Evaluating analytic gradients on quantum hardware,''
\emph{Phys. Rev. A}, vol.~99, art.~032331, 2019, doi: 10.1103/PhysRevA.99.032331.

\bibitem{qiskit2024}
A.~Javadi-Abhari \emph{et al.},
``Quantum computing with Qiskit,''
arXiv:2405.08810, 2024.
\end{thebibliography}

\end{document}
